% Source-aware proof packet for the real toric numerical-dimension theorem.
% Packet-only file.  It is not part of a live chapter.

\paragraph{Conventions and source attestations.}
Let \(M\) and \(N\) be the character and cocharacter lattices.  Lebesgue
measure on \(M_{\mathbf R}\) gives a fundamental lattice parallelepiped
volume one.  We use norm weights,
\[
 \langle m,\operatorname{Log}z\rangle=\log|\chi^m(z)|,
 \qquad
 \mathrm d^c=\frac{1}{2\pi\mathrm i}(\partial-\bar\partial).
\]
With this convention a moment polytope \(Q\) satisfies
\(\int_X\omega^n=n!\operatorname{vol}_M(Q)\).

The following four controlling sources and exact scopes are used.

\begin{enumerate}
\item The archived arXiv PDF of Junyan Cao,
\emph{A Kawamata--Viehweg--Nadel-type vanishing theorem and the
nonvanishing conjecture}, has SHA-256
\texttt{17717664b0ae2b154e76e2ecb68c9f63f09dc74be6f25bac7f04b7ee0103b0d6}.
Definition~2 on physical pages 4--5 defines a quasi-equisingular
approximation by almost positivity, analytic singularities, monotonicity,
and the asymptotic multiplier-ideal inclusion.  Definition~3 on physical
page~6 defines, for arbitrary closed positive \((1,1)\)-currents on a compact
K\"ahler manifold, the cohomological product by
\[
 \int_X\langle T_1,\ldots,T_k\rangle_{\mathrm{Cao}}\wedge u
 =\lim_j\int_X (T_{1,j})_{\mathrm{ac}}\wedge\cdots\wedge
                  (T_{k,j})_{\mathrm{ac}}\wedge u.
\]
The definition is not restricted to line-bundle classes.  Definition~4 only
specializes it to the curvature current of a singular metric on a line
bundle.  Cao's Proposition~4.4 is a line-bundle section-growth statement and
is not used here.

\item The supplied 18 July 2026 PDF
\emph{Singularities in Global Pluripotential Theory} has SHA-256
\texttt{c816eea5f06822021fa8dd75b70f82724ee47e6dcd6830729f5ecc109296b706}.
Definition~1.6.1 on printed page~31 allows a bounded remainder in a general
analytic singularity.  Definition~1.6.5 on printed page~33 says that a
gentle analytic singularity has a modification on which
\(\pi^*\varphi=c\log|g|^2+h\), with \(h\) smooth.  Theorem~1.6.2 on printed
pages 33--34 supplies a decreasing quasi-equisingular approximation all of
whose terms are gentle.  Theorem~6.2.1 on printed pages 218--221 gives
continuity of all mixed non-pluripolar masses under
\(d_{\mathcal S}\)-convergence in big classes.  Theorem~7.1.1 on printed
pages 234--236 says, in particular, that for an \(\mathcal I\)-good K\"ahler
current any quasi-equisingular approximation converges in
\(d_{\mathcal S}\).  Example~7.4.1 on printed page~256 states that every
compact-torus-invariant qpsh potential on a smooth projective toric variety
is \(\mathcal I\)-good.

\item The archived arXiv PDF of Coman--Guedj--Sahin--Zeriahi,
\emph{Toric Pluripotential Theory}, has SHA-256
\texttt{aeeff09ceabe6a62cf0879f18b8184fb1ffc4a15afbb331f50371658ecc2916d}.
Lemma~2.1 on printed page~6 gives the relation between the subgradient image
and the conjugate domain.  Lemma~2.2 on printed pages 6--7 gives, on
\((\mathbf C^*)^n\), the local complex/real Monge--Amp\`ere identity with
\(\mathrm d^c=(2\pi\mathrm i)^{-1}(\partial-\bar\partial)\), and defines its
real measure as \(n!\) times the ordinary Alexandrov measure.  Only this
local statement is used below; no global ample-class theorem from that paper
is enlarged to an arbitrary real class.

\item The archived arXiv v2 PDF of
Boucksom--Eyssidieux--Guedj--Zeriahi,
\emph{Monge--Amp\`ere equations in big cohomology classes}, has SHA-256
\texttt{43097be8459edfd948929fade34367fa7bbd0fe691d9cba1a4f27a11181214cd}.
Proposition~1.4 on printed pages 6--7 proves locality and multilinearity of
non-pluripolar products; Proposition~1.6 on printed page~7 proves that the
global product of closed positive currents on a compact K\"ahler manifold
is well defined; and the construction gives a product that puts no mass on
pluripolar sets.
\end{enumerate}

The multiplier-ideal conventions in the first two sources differ by a factor
of two.  If \(\mathcal I_{\mathrm C}\) denotes Cao's convention and
\(\mathcal I_{\mathrm B}\) the supplied book's convention, then
\[
                  \mathcal I_{\mathrm C}(a\varphi)
                  =\mathcal I_{\mathrm B}(2a\varphi).
\]
Given \(m\in\mathbf N\) and \(\delta>0\), apply the supplied book's
quasi-equisingular condition with
\(\lambda=2m\) and \(\lambda'=2m(1+\delta)\).  It gives an index \(j_0\)
for which
\[
 \mathcal I_{\mathrm C}(m(1+\delta)\varphi_{j_0})
 \subseteq\mathcal I_{\mathrm C}(m\varphi).
\]
Since the approximants decrease, the same inclusion holds for every
\(j\geq j_0\).  Thus the monotone gentle approximation is admissible in
Cao's Definition~3; no equality of the two written definitions is claimed.

\paragraph{Scope firewalls and tests.}
On a smooth projective toric variety,
\(H^{1,1}(X,\mathbf R)=N^1(X)_{\mathbf R}\); hence the new scope consists of
irrational real divisor classes, not genuinely transcendental Hodge
directions.  Multiplier ideals \(\mathcal I(c\varphi)\) still make sense for
real \(c>0\); what fails in an irrational class is a graded section algebra.

The theorem below concerns a current-level numerical dimension.  A
class-level reading is false.  On \(X=\mathbf P^1\), for
\(T=\lambda[0]\) with \(\lambda>0\), the class has numerical dimension one,
whereas Cao's absolutely continuous products discard the divisor and the
Newton body is a point.  Taking \(\lambda\notin\mathbf Q\) gives the same
counterexample in an irrational class.

For a genuinely irrational positive-dimensional test, on
\(\mathbf P^1\times\mathbf P^1\) take
\[
 T=a[\{0\}\times\mathbf P^1]
     +b\,\operatorname{pr}_2^*\omega_{\mathrm{FS}},
 \qquad a/b\notin\mathbf Q.
\]
Its Newton body is a segment.  For the product K\"ahler form with square
moment polytope, \(\operatorname{vol}(K+sQ)=bs+s^2\); the claimed formula
therefore gives current-level numerical dimension one, while the big class
has class numerical dimension two.  Purely divisorial currents give the
zero-dimensional boundary, and smooth K\"ahler currents give the
full-dimensional boundary.  The proof below also covers nef non-big and zero
classes, so bigness is deliberately not assumed.

\begin{theorem}
Let \(X=X_\Sigma\) be a smooth projective toric variety of dimension
\(n\geq1\), and let invariance mean invariance under the maximal compact
torus.  Let \(T\) be an invariant closed positive \((1,1)\)-current in a
pseudoeffective real class.  If \(F_T:N_{\mathbf R}\to\mathbf R\) is a
convex weight of \(T\) on the dense torus, define the closed Newton body
\[
                    K=\Delta(T):=
                    \overline{\operatorname{dom}F_T^*}.
\]
It is nonempty and compact and is well defined up to translation.

Let \(\omega\) be an invariant K\"ahler form, and let \(Q\) be its moment
polytope.  Normalize Lebesgue measure by the character lattice and mixed
volume by \(V(B[n])=\operatorname{vol}_M(B)\).  For every
\(0\leq k\leq n\),
\[
 \int_X\langle T^k\rangle_{\mathrm{Cao}}\wedge\omega^{n-k}
 =\int_X\langle T^k\rangle_{\mathrm{np}}\wedge\omega^{n-k}
 =n!V\bigl(K[k],Q[n-k]\bigr).
\]
Here the first product is Cao's current-level cohomological product from
quasi-equisingular approximations, and the second is the non-pluripolar
product; for \(k=0\), both products are by convention equal to \(1\).
Consequently,
\[
                    \nu_{\mathrm{Cao}}(T)
                    =\nu_{\mathrm{np}}(T)
                    =\dim K,
\]
where each numerical dimension is the largest \(k\) for which the
corresponding current-level product class is nonzero.
\end{theorem}

\begin{proof}
Choose an invariant smooth representative \(\theta\) of \(\{T\}\), average
a potential over the compact torus, and write
\[
                       T=\theta+\mathrm d\mathrm d^c\varphi.
\]
On the dense torus, a smooth invariant primitive of \(\theta\), plus
\(\varphi\), is a finite convex function
\(F=F_T:N_{\mathbf R}\to\mathbf R\).  To see that its conjugate domain is
bounded, use the finite toric affine cover.  On each cone the smooth
primitive differs from a linear function by a bounded smooth function, and
the global qpsh potential is bounded above.  Thus
\(F(x)\leq C(1+\lVert x\rVert)\).  A subgradient exists at every point
because \(F\) is finite convex on all of \(N_{\mathbf R}\), so
\(\operatorname{dom}F^*\) is nonempty; the displayed linear-growth bound
makes it bounded.  Therefore
\[
 K=\overline{\operatorname{dom}F^*}
  =\overline{\bigcup_x\partial F(x)}
\tag{1}
\]
is nonempty and compact.  Changing the invariant primitive by an affine
function translates \(K\), proving the asserted intrinsicness.  The closure
in (1) is essential: for
\(F(x)=\log(1+\mathrm e^x)\), the gradient image is \((0,1)\), while the
conjugate domain is \([0,1]\).

Fix the decreasing gentle quasi-equisingular approximation
\((\varphi_j)\) supplied by the cited regularization theorem, and put
\[
                  R_j=\theta+\mathrm d\mathrm d^c\varphi_j.
\]
After comparing its auxiliary K\"ahler form with \(\omega\), write
\(R_j\geq-\varepsilon_j\omega\), where
\(\varepsilon_j\downarrow0\).  For \(s>0\), set
\(S_s=T+s\omega\).  Deleting finitely many terms makes the same sequence a
quasi-equisingular approximation for \(S_s\), because
\(R_j+s\omega\geq0\), while addition of a smooth form changes neither
analytic singularities nor multiplier ideals.  Since
\[
 (R_j+s\omega)_{\mathrm{ac}}=(R_j)_{\mathrm{ac}}+s\omega,
\]
Cao's Definition~3 and the ordinary binomial formula give
\[
 \int_X\langle S_s^n\rangle_{\mathrm{Cao}}
 =\sum_{k=0}^n\binom nk s^{n-k}
   \int_X\langle T^k\rangle_{\mathrm{Cao}}
                  \wedge\omega^{n-k}.
\tag{2}
\]
This is an identity for one fixed approximation, not a class-varying
continuity assertion.

We next identify the left side of (2).  The current \(S_s\) dominates
\(s\omega\), so it is K\"ahler.  Its potential is toric and hence
\(\mathcal I\)-good by the cited Example~7.4.1.  The cited Theorem~7.1.1
therefore gives \(d_{\mathcal S}\)-convergence of the fixed
quasi-equisingular approximation to \(S_s\), and Theorem~6.2.1 gives
convergence of its top non-pluripolar masses.

The use of gentle singularities is load bearing.  For fixed \(s\) and each
term in the positive tail, Definition~1.6.5 supplies a modification
\(\pi_j:Y_j\to X\) and an exact decomposition
\[
                 \pi_j^*(R_j+s\omega)=[E_j]+\beta_{j,s},
\tag{3}
\]
where \(E_j\) is an effective real divisor and \(\beta_{j,s}\) is smooth.
Off \(\operatorname{Supp}E_j\), the latter is the pullback of a positive
current, hence is semipositive there and, by continuity, everywhere.  The
singular and exceptional analytic sets have Lebesgue measure zero, while
the non-pluripolar product ignores them.  On their complement both products
are the ordinary smooth wedge.  Thus
\[
 \int_X(R_j+s\omega)_{\mathrm{ac}}^n
 =\int_{Y_j}\beta_{j,s}^n
 =\int_X\langle(R_j+s\omega)^n\rangle_{\mathrm{np}}.
\tag{4}
\]
The smooth residual form in (3) would not follow from a merely bounded
remainder in the general definition of analytic singularities.  Taking the
limit in (4) and using the preceding \(d_{\mathcal S}\)-continuity yields
\[
 \int_X\langle S_s^n\rangle_{\mathrm{Cao}}
 =\int_X\langle S_s^n\rangle_{\mathrm{np}}.
\tag{5}
\]

Let \(G\) be the smooth strictly convex weight of \(\omega\) on the dense
torus, so that
\(\overline{\operatorname{dom}G^*}=Q\).  The weight of \(S_s\) is
\(F_s=F+sG\).  It is finite and has a linear upper-growth bound, so its
conjugate domain is again nonempty and bounded.  We claim
\[
             \overline{\operatorname{dom}F_s^*}=K+sQ.
\tag{6}
\]
Here is a closure-safe proof.  If \(H\) is a finite convex function with
bounded conjugate domain \(B=\overline{\operatorname{dom}H^*}\), its
recession function is \(H^\infty=h_B\).  Indeed, Fenchel's inequality gives
\(H^\infty(v)\geq\langle p,v\rangle\) for every
\(p\in\operatorname{dom}H^*\); conversely, integration of one-dimensional
directional derivatives, all of whose subgradients lie in \(B\), gives
\(H(x+tv)-H(x)\leq t h_B(v)\).  Recession functions add, and support
functions determine compact convex bodies.  Therefore
\[
 F_s^\infty=F^\infty+sG^\infty
            =h_K+s h_Q=h_{K+sQ},
\]
which proves (6).  This avoids the generally false unclosed-domain identity
for conjugates of sums.

On \((\mathbf C^*)^n\), the finite convex function \(F_s\) is locally
bounded, so the local Bedford--Taylor product is defined.  The cited local
complex/real formula identifies its logarithmic pushforward with the real
Monge--Amp\`ere measure, whose total mass is \(n!\) times the Lebesgue
measure of the subgradient image.  The interior of
\(\operatorname{dom}F_s^*\) lies in that image by the cited Lemma~2.1, and
the image lies in the domain; since (6) is full dimensional, passing to its
closure changes only a Lebesgue-null boundary.  The complement of the dense
torus in \(X\) is analytic and hence pluripolar, so the non-pluripolar measure
does not charge it.  Consequently,
\[
 \int_X\langle S_s^n\rangle_{\mathrm{np}}
       =n!\operatorname{vol}_M(K+sQ).
\tag{7}
\]
This derives (7) from a local statement and the pluripolar-boundary
argument; it does not silently extend a published global ample-class
statement.

Minkowski's polynomial formula is
\[
 \operatorname{vol}_M(K+sQ)
 =\sum_{k=0}^n\binom nk
       V(K[k],Q[n-k])s^{n-k}.
\tag{8}
\]
Equations (2), (5), (7), and (8) hold for every \(s>0\).  Equality of these
polynomials gives
\[
 \int_X\langle T^k\rangle_{\mathrm{Cao}}\wedge\omega^{n-k}
       =n!V(K[k],Q[n-k]).
\tag{9}
\]
By the cited multilinearity of non-pluripolar products, expanding
\(\langle(T+s\omega)^n\rangle_{\mathrm{np}}\) and comparing it with
(7)--(8) similarly gives
\[
 \int_X\langle T^k\rangle_{\mathrm{np}}\wedge\omega^{n-k}
       =n!V(K[k],Q[n-k]).
\tag{10}
\]

It remains to determine which coefficients vanish.  Put \(d=\dim K\).
For \(k>d\), translate \(K\) into a \(d\)-dimensional linear subspace and
place it in a bounded box there.  The volume of
\((t_1+\cdots+t_k)K+sQ\), as a polynomial in the \(t_i\), has degree at most
\(d\) in \(t_1+\cdots+t_k\); hence its \(t_1\cdots t_k\) coefficient is zero.
This coefficient is a positive universal multiple of
\(V(K[k],Q[n-k])\), so that mixed volume vanishes.  For \(1\leq k\leq d\),
translate \(K\) and choose \(k\) linearly independent segments in it.
Complete their directions by \(n-k\) segments contained, after translation
and scaling, in the interior of \(Q\).  The mixed volume of these segments
is a positive multiple of the absolute determinant of their directions, and
monotonicity of mixed volume gives
\[
 V(K[k],Q[n-k])>0.
\tag{11}
\]
For \(k=0\), positivity is simply \(V(Q[n])=\operatorname{vol}_M(Q)>0\).
Cao's product class is pseudoeffective, and the non-pluripolar product is
represented by a positive current.  Such a class is nonzero exactly when
its pairing with the strictly positive form \(\omega^{n-k}\) is positive.
Equations (9)--(11) therefore give
\[
                    \nu_{\mathrm{Cao}}(T)
                    =\nu_{\mathrm{np}}(T)
                    =d.
\]
\end{proof}
